\section{\texorpdfstring{$\Lzm$}{L0}: Axiomatic Semantics}
\label{sec:l0}

% ===========================================================================
\subsection{Histories}
\label{sec:l0-histories}

Everything here is a function of the \gls{history} (\cref{def:history}): the
\glspl{operation} issued, ordered per rank. No timing, message order or
interleaving is an input.

\begin{notationx}\label{not:history}
Throughout, $\Hist$ is a history over a rank set $\Ranks$ with a distance
\gls{comm-graph} $\cgraph$ (\cref{def:comm-graph}). We write $\Hist_{\opput}$ and
$\Hist_{\opget}$ for the \glspl{put} and \glspl{get} in $\Hist$, and
$\Hist|_p \defeq \{o \in \Hist : \opsrc{o} = p\}$ for the operations issued by
process $\proc{p}$, totally ordered by $\po$.
\end{notationx}

% ===========================================================================
\subsection{Restrictions}
\label{sec:l0-restrictions}

A history is not automatically meaningful. DSTM buys its simple memory model
by pushing the difficulty onto the program, as eight restrictions in three
groups:



\begin{description}[leftmargin=1.1cm,style=nextline,font=\normalfont\bfseries]
  \item[Semantic] \cref{res:write-disjoint} alone makes the store well
    defined. Without it ``the'' value of a \gls{cell} does not exist, and
    there is nothing to be correct \emph{about}.
  \item[Mechanical] \cref{res:monotone,res:lookahead,res:staleness,res:phase}
    make the store \emph{computable} by the mechanism of \cref{sec:l1}. A
    program violating them has a well-defined meaning DSTM cannot deliver.
  \item[Resource and lifecycle] \cref{res:horizon,res:collective,%
    res:terminate} bound memory and govern \gls{session} structure.
\end{description}

None of the eight is currently checked by the implementation; see
\crefrange{dev:zero-distance}{dev:barrier-doc}. \Cref{ch:restrictions}
collects them with proposed runtime checks.

\subsubsection*{Mechanical restrictions}

\begin{restr}[Monotone origin time]\label{res:monotone}
For every process $\proc{p}$ and every $o_1, o_2 \in \Hist|_p$ with $o_1 \po o_2$,
\[
  \optime{o_1} \le \optime{o_2}.
\]
Equivalently: a process never moves backwards along the timeline.
\devmark{dev:monotone-origin}
\end{restr}

\begin{restr}[Lookahead]\label{res:lookahead}
For every put $o \in \Hist_{\opput}$ with $p = \opsrc{o}$ and
$q = \optgt{o}$,
\[
  \opdelta{o} \;\ge\; \lag{p}{q},
  \qquad\text{and}\qquad
  q \in \nbr{p} \ \text{ or }\ q = p .
\]
The local case degenerates to $\opdelta{o} \ge 0$.
\devmark{dev:lookahead-uncheck}
\end{restr}

\begin{restr}[Staleness bound]\label{res:staleness}
For every get $o \in \Hist_{\opget}$ with $p = \opsrc{o}$ and
$q = \optgt{o}$,
\[
  -\opdelta{o} \;\le\; \lag{p}{q},
  \qquad\text{and}\qquad
  q \in \nbr{p} \ \text{ or }\ q = p .
\]
The local case degenerates to $\opdelta{o} \le 0$ combined with
$-\opdelta{o} \le 0$, hence $\opdelta{o} = 0$: a process reading its own \gls{region}
may only read its current effective time.
\devmark{dev:lookahead-uncheck}
\end{restr}

\Cref{res:lookahead,res:staleness} are duals, and reading them together is the
key to the whole model:

\begin{center}
\emph{The \gls{lag} is a lower bound on how long a write takes to
become visible,\\ and an upper bound on how far back a read may look.}
\end{center}

A write may take longer than $\lag{p}{q}$ and a read may look back less far;
what neither may do is beat the lag. Everything in \cref{sec:l1} is a
consequence.

\begin{restr}[Phase order]\label{res:phase}
For every process $\proc{p}$ and every $w \in \Hist_{\opput} \cap \Hist|_p$,
$g \in \Hist_{\opget} \cap \Hist|_p$ with
$\optgt{w} = \optgt{g}$ and $\opeff{w} = \opeff{g}$,
\[
  w \po g .
\]
That is: within one process, every write to a given version is issued before any
read of that version. \devmark{dev:phase-strictness}
\end{restr}

\Cref{res:phase} is stated per rank, which looks too weak: surely a
\emph{remote} write arriving after a read is served is equally fatal? It is,
but cannot happen.

\begin{proposition}[Minimality of the phase restriction]
\label{prop:phase-minimality}
Let $\Hist$ satisfy \cref{res:monotone,res:lookahead,res:phase}, and let
$g$ be a get targeting $(q,\eff)$ and $w$ a put targeting $(q,\eff)$ with
$\opsrc{w} \ne \opsrc{g}$. If $\opsrc{w} \ne q$ then
$\optime{w} \le \eff - \lag{\opsrc{w}}{q} < \eff$, whereas any operation
issued by $\opsrc{w}$ after a read of $(q,\eff)$ could be served has origin
time at least $\eff$. Hence the cross-process case of the phase order is implied,
and the per-rank statement of \cref{res:phase} is the only genuine
restriction.
\end{proposition}

\begin{proof}
By \cref{res:lookahead}, $\opdelta{w} \ge \lag{\opsrc{w}}{q} \ge 1$ using
(G3) and $\opsrc{w} \ne q$. Since $\opeff{w} = \optime{w} + \opdelta{w} =
\eff$, we get $\optime{w} \le \eff - 1 < \eff$. By \cref{res:monotone} the
\glspl{origin-time} of $\opsrc{w}$ are nondecreasing, so every put by $\opsrc{w}$
targeting $(q,\eff)$ occurs while that process's clock is still at most
$\eff - 1$. The publish guard of \cref{sec:l1-rules} does not fire until every
neighbour's \gls{known-time} $\kn{q}{r}$ satisfies $\kn{q}{r} + \lag{q}{r} > \eff$,
which for $r = \opsrc{w}$ requires $\opsrc{w}$ to have advanced strictly past
$\eff - \lag{q}{\opsrc{w}}$, i.e.\ past the last origin time at which it
could have produced such a put.

The remaining case is $\opsrc{w} = q \ne \opsrc{g}$, a local put on the target
process concurrent with a remote read. By \cref{res:staleness} applied to $g$,
$\eff \le \optime{g} - \lag{\opsrc{g}}{q} < \optime{g}$, and the read cannot
be served before $q$'s own clock reaches $\eff$; a local put at $\eff$ issued
after that point would violate \cref{res:monotone} unless it has
$\opdelta{} = 0$ and origin exactly $\eff$, which is the case
\cref{res:phase} governs on process $\proc{q}$ itself.
\end{proof}

This explains a puzzling feature of the implementation: the publish guard
tests $\clk{p} \ge \eff$, not $\clk{p} > \eff$
\devmark{dev:phase-strictness}. The non-strict test leaves exactly one hole
--- a local, zero-\gls{delta} write at the process's current time --- and
\cref{res:phase} plugs that hole and nothing more. A strict guard would close
it automatically but forbid a process from reading its own current tick, which
is the entire \glslink{push-model}{push idiom} (\cref{sec:exec-idioms}).

\subsubsection*{Semantic restriction}

\begin{restr}[Write disjointness]\label{res:write-disjoint}
For all distinct $w_1, w_2 \in \Hist_{\opput}$,
\[
  \fpw{w_1} \cap \fpw{w_2} = \emptyset .
\]
Equivalently: every cell is written at most once in the entire execution.
\devmark{dev:write-disjoint}
\end{restr}

Because a \gls{cell} carries an \gls{effective-time}, this is weaker than it
sounds: two writes to one address at different ticks do not conflict, only two
at the same tick. The \glslink{mailbox}{mailbox idiom}
(\cref{sec:gloss-idioms}) discharges it structurally.

\subsubsection*{Resource and lifecycle restrictions}

\begin{restr}[Bounded horizon]\label{res:horizon}
There is a constant $\delta_{\max}$ such that $\opdelta{o} \le \delta_{\max}$
for every $o \in \Hist_{\opput}$.
\end{restr}

This affects no result about values; only \cref{thm:space-bound} needs it. A
put with unbounded \gls{delta} creates a version that stays writable
indefinitely, and memory grows without limit.

\begin{restr}[Collective agreement]\label{res:collective}
Every collective entry point --- session start and end, the \gls{termination}
\gls{barrier}, and the creation of any derived structure --- is invoked by every
process, the same number of times, in the same order, with matching arguments.
\devmark{dev:noncollective}
\end{restr}

\begin{restr}[Termination bound]\label{res:terminate}
The argument $\ell$ passed to the terminating barrier satisfies
\[
  \ell \;\ge\; \max_{o \in \Hist} \opeff{o},
\]
and every process passes the same $\ell$. \devmark{dev:clock-conflation}
\end{restr}

\begin{definition}[Legality]\label{def:legal}
$\Legal{\Hist}$ holds when $\Hist$ satisfies
\crefrange{res:monotone}{res:terminate}.
\end{definition}

% ===========================================================================
\subsection{The store}
\label{sec:l0-store}

With legality in hand the semantics is a single equation.

\begin{definition}[Denotation of a history]\label{def:store-h}
For a history $\Hist$ satisfying \cref{res:write-disjoint}, the store
$\StoreH : \Cells \to \Bytes$ is
\[
  \StoreH\bigl(\cell{q}{\eff}{a}\bigr) \;\defeq\;
  \begin{cases}
    v_j
      & \text{if } \exists\, w = \Put{p}{q}{t}{\delta}{a_0}{v} \in \Hist
        \text{ with } t + \delta = \eff \\[-1pt]
      & \text{and } j = a - a_0 \in \{0, \dots, |v| - 1\}, \\[4pt]
    0 & \text{otherwise.}
  \end{cases}
\]
\end{definition}

\begin{lemma}[Well-definedness]\label{lem:store-welldef}
If $\Hist$ satisfies \cref{res:write-disjoint} then $\StoreH$ is a
well-defined total function.
\end{lemma}

\begin{proof}
Totality is immediate from the second clause. For single-valuedness, suppose
two puts $w_1 \ne w_2$ both satisfy the first clause for the cell
$c = \cell{q}{\eff}{a}$. Then $c \in \fpw{w_1}$ and $c \in \fpw{w_2}$ by
\cref{def:footprint}, contradicting \cref{res:write-disjoint}.
\end{proof}

\Cref{def:store-h} is the whole of the DSTM memory semantics: no clocks, no
lags, no \glspl{version}, no statuses, no messages. Those exist in
\cref{sec:l1} solely to compute this function without a global \gls{barrier}.

\begin{definition}[Read result]\label{def:read-result}
A get $g = \Get{p}{q}{t}{\delta}{a}{\ell}$ \emph{returns}
\[
  \StoreH(\cell{q}{t+\delta}{a})
  \;\cdots\;
  \StoreH(\cell{q}{t+\delta}{a+\ell-1}).
\]
\end{definition}

\begin{definition}[Default-valued read]\label{def:get-or-default}
A default-valued get carries an extra argument $d \in \Values$ with $|d| =
\ell$ and returns $d$ when $t + \delta < 0$, and behaves as
\cref{def:read-result} otherwise. This is the only construct that consults the
sign of an effective time; it exists so that the first ticks of a simulation,
whose producers have no history yet, need not be special-cased by the
application.
\end{definition}

\begin{definition}[Reads-from]\label{def:reads-from}
$w \rf g$ holds when $w \in \Hist_{\opput}$, $g \in \Hist_{\opget}$ and
$\fpw{w} \cap \fpr{g} \ne \emptyset$. By \cref{res:write-disjoint} each byte
returned by $g$ is supplied by at most one such $w$; bytes with no such $w$
read as $0$.
\end{definition}

% ===========================================================================
\subsection{Determinacy}
\label{sec:l0-determinacy}

$\StoreH$ depends on $\Hist$ alone, so fixing the \gls{history} fixes every
read. The question is whether the history itself is fixed: control flow
depends on what is read, which depends on what others wrote. The restrictions
make that circle well founded.

\begin{definition}[Deterministic program]\label{def:det-program}
A program is \emph{deterministic} when the operation each process issues next,
and the values it writes, are functions of that process's inputs and of the
values it has read so far.
\end{definition}

\begin{theorem}[Determinacy]\label{thm:determinism}
Let $P$ be a deterministic program all of whose executions are legal
(\cref{def:legal}). Then $P$ has a unique history $\Hist(P)$ and a unique
store $\Sigma_{\Hist(P)}$. In particular, every read returns the same value in
every execution, independently of scheduling, message order, and the relative
progress of processes.
\end{theorem}

\begin{proof}
Order operations by the lexicographic process
$\rho(o) \defeq \bigl(\opeff{o},\, \iota(o)\bigr)$, where $\iota(o)$ is the
position of $o$ in the \gls{program-order} of $\opsrc{o}$. We show that the value
returned by a get $g$ depends only on operations of strictly smaller $\rho$,
so that $\Hist(P)$ is determined by induction.

Let $g$ be a get targeting $(q,\eff)$ and let $w$ be any put with
$\fpw{w} \cap \fpr{g} \ne \emptyset$; then $\optgt{w} = q$ and
$\opeff{w} = \eff$, so the first components agree. Two cases.

\emph{Case $\opsrc{w} \ne \opsrc{g}$.} By \cref{prop:phase-minimality}, $w$ is
issued by a process whose clock is strictly below $\eff$ at the time, while $g$'s
issuer has clock at least $\eff$; the two are ordered by the causal chain of
\glspl{heartbeat}, and $w$'s issue necessarily precedes $g$'s completion. Its value
therefore depends only on reads that $\opsrc{w}$ performed at origin times
$\le \eff - \lag{\opsrc{w}}{q} < \eff$, all of strictly smaller first
component by (G3).

\emph{Case $\opsrc{w} = \opsrc{g} = p$.} Then $\optgt{w} = q$ and
$\opeff{w} = \opeff{g}$, so \cref{res:phase} gives $w \po g$, hence
$\iota(w) < \iota(g)$ and $\rho(w) < \rho(g)$.

In both cases every put contributing to $g$'s result has strictly smaller
$\rho$. The relation is well founded because \glspl{effective-time} are bounded below
by $0$ on any version that exists and program-order positions are natural
numbers. Determinism of $P$ then fixes each successive operation, and
\cref{lem:store-welldef} fixes each read result.
\end{proof}

\begin{corollary}[Equivalence with lockstep execution]\label{thm:bsp}
Let $P$ be as in \cref{thm:determinism}, and let $\mathrm{BSP}(P)$ be the
execution in which every process completes tick $t$ and passes a global barrier
before any process begins tick $t+1$. Then $\mathrm{BSP}(P)$ is a legal execution
of $P$ and produces $\Sigma_{\Hist(P)}$. Consequently every DSTM execution of
$P$ returns exactly the values the barrier-synchronised execution would.
\end{corollary}

\begin{proof}
$\mathrm{BSP}(P)$ satisfies every restriction, since a global barrier at each
tick boundary only adds ordering. By \cref{thm:determinism} the history is
unique, so $\mathrm{BSP}(P)$ realises $\Hist(P)$.
\end{proof}

\Cref{thm:bsp} licenses the whole design: removing the per-tick \gls{barrier}
and letting \glspl{process} drift changes \emph{when} work happens and
\emph{how much memory} is live, provably not \emph{what is computed}. Unlike
optimistic simulation nothing is ever discarded, because nothing is ever
wrong.

% ===========================================================================
\subsection{Objects as transition systems}
\label{sec:l0-ts}

The \gls{store} is static; what moves are the clocks, known times and
\glspl{version-status}. Each is a transition system, and \cref{sec:l1}'s
global system is their composition. Separating them shows which invariants are
local and which belong to the composition.

\begin{definition}[Transition system]\label{def:ts}
A \emph{transition system} is a tuple
$\TS{} = \langle \States, \Init, \Labels, \steps \rangle$ with $\States$ a set
of states, $\Init \in \States$ the initial state, $\Labels$ a set of labels
and ${\steps} \subseteq \States \times \Labels \times \States$. We write
$s \stepl{\lambda} s'$ for $(s,\lambda,s') \in {\steps}$.
\end{definition}

\begin{definition}[Clock]\label{ts:clock}
$\TS{\mathrm{clk}}(p) \defeq \langle \Time,\; 0,\;
 \{\optick(t) : t \in \Time\},\; \steps \rangle$ with
\[
  c \stepl{\optick(t)} t \quad\text{iff}\quad t > c .
\]
Monotone and deterministic in the label.
\end{definition}

\begin{definition}[Known time]\label{ts:known}
$\TS{\mathrm{kn}}(p,q) \defeq \langle \Time,\; 0,\;
 \{\mathsf{hear}(t) : t \in \Time\},\; \steps \rangle$ with
\[
  k \stepl{\mathsf{hear}(t)} \max(k, t).
\]
Total --- every label is enabled in every state --- and monotone.
\end{definition}

\begin{definition}[Version status]\label{ts:version}
$\TS{\mathrm{st}}(p,\eff) \defeq \langle \Status,\; \Absent,\;
 \{\mathsf{create}, \mathsf{publish}, \mathsf{retire}, \mathsf{evict}\},\;
 \steps \rangle$ with
\[
  \Absent \stepl{\mathsf{create}} \Wr
  \stepl{\mathsf{publish}} \Rd
  \stepl{\mathsf{retire}} \Dc
  \stepl{\mathsf{evict}} \Absent .
\]
Every transition is monotone in the order of \cref{def:status} except
$\mathsf{evict}$, which returns the system to $\Absent$. That single
non-monotone edge is why \cref{thm:discard-safety} is needed: it must be
impossible to distinguish an evicted version from one that never existed,
which holds exactly when nothing will ever ask for it again.
\end{definition}

\begin{definition}[Cell]\label{ts:cell}
$\TS{\mathrm{cell}}(c) \defeq \langle \Bytes,\; 0,\;
 \{\mathsf{write}(v) : v \in \Bytes\},\; \steps \rangle$ with
$b \stepl{\mathsf{write}(v)} v$, subject to the guard that
$\mathsf{write}$ is enabled only while the enclosing version's status is
$\Wr$. Under \cref{res:write-disjoint} at most one $\mathsf{write}$ ever
occurs, so the reachable states are exactly $\{0\} \cup \{v\}$ and the
system is confluent.
\end{definition}

\begin{definition}[Allocation]\label{ts:alloc}
$\TS{\mathrm{alloc}}(p) \defeq \langle \States,\; \emptyset,\;
 \{\mathsf{alloc}(n,h,I), \mathsf{free}(h)\},\; \steps \rangle$, where
$\States$ is the set of well-formed allocation states
(\cref{def:allocation}) and
\[
  \allocof{p} \stepl{\mathsf{alloc}(n,h,I)} \allocof{p}[\,I \mapsto h\,],
  \qquad
  \allocof{p} \stepl{\mathsf{free}(h)} \restrict{\allocof{p}}{\Addrof{p}
    \setminus \extent{p}{h}},
\]
the first enabled when $h$ is fresh and $I$ is a length-$n$ interval disjoint
from $\live{p}$. The label carries $I$: the system is \emph{nondeterministic},
and every resolution of that nondeterminism is a conforming allocator
(\cref{rem:alloc-policy}).
\end{definition}

\begin{remark}[Where the difficulty lives]
Each of \crefrange{ts:clock}{ts:alloc} is trivial in isolation: two monotone
counters, a four-state chain, a single-assignment register and a disjointness
bookkeeper. All the difficulty is in the \emph{guards} coupling them --- when
$\mathsf{publish}$ and $\mathsf{retire}$ may fire, which is a question about
other processes. \Cref{sec:l1} supplies them.
\end{remark}

\Cref{ts:alloc} is the one object with no such guard: allocation is purely
local, involves no message and constrains no other process. That is why it
does not appear in \cref{sec:l1-rules} --- its work is done before any
operation is issued.

% ===========================================================================
\subsection{The correctness criterion}
\label{sec:l0-correctness}

\begin{definition}[Correct implementation]\label{def:correctness}
An implementation is \emph{correct} if, for every deterministic program $P$
all of whose executions are legal, every execution of $P$
\begin{enumerate}[label=(\roman*),leftmargin=1.4cm]
  \item terminates, with every get having completed, and
  \item returns, for each get, the bytes prescribed by
    \cref{def:read-result} with respect to $\Sigma_{\Hist(P)}$.
\end{enumerate}
\end{definition}

Clause (ii) is safety and is discharged by \cref{thm:read-soundness}; clause
(i) is liveness and is discharged by \cref{thm:deadlock-freedom} together with
\cref{thm:barrier-correct}. Both are proved against $\Lom$ in the next
section.
